\section{Besov球的覆盖数}
\label{sec:besov-covering}

前两节回答的是“给一个函数，用多少项能达到给定误差”。另一类问题是：为整类函数预先准备有限个代表，使类中每个函数都离某个代表足够近，至少需要多少个代表？第~\ref{chap:sobolev-spaces}章的紧嵌入只说明Sobolev球的覆盖数有限，没有给出增长阶。本节回到第~\ref{sec:besov-scale-truncation}节的$B^s_{2,q}$有界类，在$L^2$距离下算出覆盖数随精度的增长：固定尺度截断控制尾部，逐尺度覆盖保留的系数给出上界，再用只在单一尺度上变化的子类给出同阶的下界。

\subsection{函数类、距离与精度}
\label{sec:besov-cover-setting}

固定单位周期、$s>0$、$1\leq q\leq\infty$与$R>0$，令
\[
\mathcal F_R=\{f\in L^2(\mathbb T):\lVert f\rVert_{B^s_{2,q}}\leq R\}.
\]
沿用第~\ref{chap:metric-spaces}章的覆盖数：记$N(\varepsilon,\mathcal F_R)$为覆盖$\mathcal F_R$所需的半径为$\varepsilon$的$L^2$闭球的最少个数，球心可以取在整个$L^2(\mathbb T)$中。若要求球心属于$\mathcal F_R$，由那里的比较$N_X(\varepsilon,A)\leq N_A(\varepsilon,A)\leq N_X(\varepsilon/2,A)$，只需把半径调整一个固定的因子，下面的增长阶不变。

我们关心$\varepsilon\to0$时$\log N(\varepsilon,\mathcal F_R)$怎样增长，即度量熵。$\mathcal F_R$在$L^2$中全有界：统一截断界说明它离有限维子空间$V_J$一致地近，而它在$V_J$上的投影是有界集。下面把这个定性论证做成定量计数。

\subsection{逐尺度精度分配}
\label{sec:besov-precision-allocation}

由第~\ref{sec:besov-uniform-truncation}小节，类中每个函数满足$E_J(f)\leq C_sR2^{-Js}$，并且由式\eqref{eq:besov-wavelet-norm-l2}，每层系数满足$\lVert a\rVert_{\ell^2}\leq C_0R$，$\lVert c_j\rVert_{\ell^2}\leq C_0R2^{-js}$。覆盖的构造分两部分：尾部直接舍去，保留的系数逐层覆盖。

\emph{尾部。}设$0<\varepsilon<2C_sR\,2^{-(j_0+1)s}$，取$J$为使$C_sR2^{-Js}\leq\varepsilon/2$的最小整数。这一假设保证$J>j_0+1$，并且由最小性$C_sR2^{-(J-1)s}>\varepsilon/2$。于是每个函数的尾部误差$E_J(f)\leq\varepsilon/2$，并且
\begin{equation}
\frac{R}{\varepsilon}\leq\frac{2^{Js}}{2C_s},\qquad2^J<2\Bigl(\frac{2C_sR}{\varepsilon}\Bigr)^{1/s}.
\label{eq:besov-cover-scale}
\end{equation}

\emph{保留的系数。}把粗项$a$并入起始一层，组成维数为$2^{j_0}+2^{j_0}=2^{j_0+1}$的块$(a,c_{j_0})$，半径至多$\sqrt2C_0R$；其余每个$j_0<j<J$的块$c_j$维数为$2^j$，半径至多$C_0R2^{-js}$。对第$j$块分配覆盖精度
\[
\delta_j=\frac\varepsilon2\,2^{-(J-j)/2},\qquad j_0\leq j<J .
\]
越粗的层分到越大的精度，因为它们半径大而维数小；越接近截断尺度的层精度越细。总和满足
\[
\sum_{j_0\leq j<J}\delta_j^2=\frac{\varepsilon^2}{4}\sum_{j_0\leq j<J}2^{-(J-j)}<\frac{\varepsilon^2}{4}\sum_{m\geq1}2^{-m}=\frac{\varepsilon^2}4 .
\]

\emph{合并。}对每块取一个$\delta_j$覆盖，从每块的覆盖中各选一个中心，组成$V_J$中的一个函数$g$。给定$f\in\mathcal F_R$，在每块中选取与$f$的该块系数相距不超过$\delta_j$的中心。由正交性，误差的平方等于各块误差平方与尾部误差平方之和：
\[
\lVert f-g\rVert_2^2\leq\sum_{j_0\leq j<J}\delta_j^2+E_J(f)^2<\frac{\varepsilon^2}4+\frac{\varepsilon^2}4<\varepsilon^2 .
\]
所以这些组合中心构成$\mathcal F_R$的$\varepsilon$覆盖，个数是各块覆盖个数的乘积。

\subsection{对数覆盖数的上界}
\label{sec:besov-entropy-upper}

现在计数。每块通过正交坐标等距地对应于欧氏空间中的球。由第~\ref{chap:metric-spaces}章，$n$维实赋范空间中半径为$\rho$的闭球可以用不超过$(1+2\rho/\delta)^n$个半径为$\delta$的闭球覆盖。

\begin{theorem}[对数覆盖数的上界]
\label{thm:besov-entropy-upper}
设$s>0$，$1\leq q\leq\infty$。存在只依赖$s,q$与所选小波的常数$c,C>0$，使对$0<\varepsilon\leq cR$，
\[
\log N(\varepsilon,\mathcal F_R)\leq C\Bigl(\frac R\varepsilon\Bigr)^{1/s}.
\]
\end{theorem}

\begin{proof}
取$c=C_s2^{-(j_0+1)s}$，则$\varepsilon\leq cR$保证上一小节的假设成立。记$\rho_j$为第$j$块的半径。对$j_0<j<J$，由式\eqref{eq:besov-cover-scale}与$C_0\leq C_s$，
\[
\frac{2\rho_j}{\delta_j}\leq\frac{4C_0R\,2^{-js}2^{(J-j)/2}}{\varepsilon}\leq\frac{2C_0}{C_s}2^{(J-j)s}2^{(J-j)/2}\leq2\cdot2^{(J-j)(s+1/2)},
\]
所以$\log(1+2\rho_j/\delta_j)\leq\log3+(J-j)(s+1/2)\log2\leq C_1(1+J-j)$。起始块的半径多一个因子$\sqrt2\,2^{j_0s}$，只改变常数。按块相乘再取对数，
\[
\log N(\varepsilon,\mathcal F_R)\leq\sum_{j_0\leq j<J}2^{j+1}C_1(1+J-j)\leq2C_1\,2^J\sum_{m\geq1}2^{-m}(1+m)=6C_1\,2^J ,
\]
这里令$m=J-j$，并用了$\sum_{m\geq1}2^{-m}=1$与$\sum_{m\geq1}m2^{-m}=2$。再由式\eqref{eq:besov-cover-scale}，$2^J<2(2C_sR/\varepsilon)^{1/s}$，即得结论。
\end{proof}

估计的对象是对数覆盖数，不是覆盖数本身：覆盖数按$\exp(C(R/\varepsilon)^{1/s})$增长。计数的主要部分来自接近截断尺度的几层，那里维数最大；每层贡献“维数乘以$1+J-j$”，几何权重使总和与最细保留层的维数$2^J$同阶。常数依赖固定的$s,q$与小波，适用范围是$0<\varepsilon\leq cR$；$\varepsilon$与$R$同阶或更大时，少数几个代表就够了。

\subsection{覆盖数的下界}
\label{sec:besov-entropy-lower}

上界是否太粗？我们用一个只在单一尺度上变化的子类核对。

\begin{theorem}[对数覆盖数的下界]
\label{thm:besov-entropy-lower}
在定理~\ref{thm:besov-entropy-upper}的条件下，存在常数$c',c''>0$，使对$0<\varepsilon\leq c''R$，
\[
\log N(\varepsilon,\mathcal F_R)\geq c'\Bigl(\frac R\varepsilon\Bigr)^{1/s}.
\]
\end{theorem}

\begin{proof}
令$\alpha=1/C_0$。对$j\geq j_0$，考虑只允许第$j$层系数变化、其余系数为零的函数$f=\sum_kx_k\psi_{j,k}$，$\lVert x\rVert_{\ell^2}\leq\alpha R2^{-js}$。它的系数范数为$2^{js}\lVert x\rVert_{\ell^2}\leq\alpha R$，所以$\lVert f\rVert_{B^s_{2,q}}\leq C_0\alpha R=R$，这些函数组成$\mathcal F_R$的子集$\mathcal G_j$。$\mathcal G_j$等距于$2^j$维欧氏空间中半径$\rho=\alpha R2^{-js}$的闭球。

覆盖$\mathcal F_R$的球也覆盖$\mathcal G_j$。到第$j$层子空间的正交投影不增大与该子空间中点的距离，所以可以把球心都投影到这个子空间中。$N$个半径$\varepsilon$的球覆盖半径$\rho$的球，比较体积得$N\varepsilon^{2^j}\geq\rho^{2^j}$，即$\log N\geq2^j\log(\rho/\varepsilon)$。

选择尺度：数列$\alpha R2^{-js}$（$j\geq j_0$）每步缩小因子$2^s$，若$\varepsilon\leq\alpha R2^{-j_0s}/4$，则恰有一个$j$满足$4\varepsilon\leq\alpha R2^{-js}<2^{s+2}\varepsilon$。对这个$j$，$\rho/\varepsilon\geq4$，所以$\log N\geq2^j\log4$；又$2^{js}>\alpha R/(2^{s+2}\varepsilon)$，即$2^j>(\alpha/2^{s+2})^{1/s}(R/\varepsilon)^{1/s}$。合并即得结论，$c'=\log4\cdot(\alpha/2^{s+2})^{1/s}$。
\end{proof}

上下界合起来，$\log N(\varepsilon,\mathcal F_R)$与$(R/\varepsilon)^{1/s}$同阶，增长阶由光滑阶数$s$决定，与汇总指数$q$无关：$q$改变球的大小，但单尺度子类与逐尺度上界只用到每层的统一界，两端的量级相同。这个结论依赖本节固定的设定，即单位周期、$p=2$的Besov球与$L^2$距离；其他积分指数、误差范数或区域需要另行分析，例如目标范数与$p$不同时，单尺度子类不再是欧氏球。与Sobolev紧嵌入只给出“有限”相比，这里的小波坐标把函数类分解为一串维数与半径都已知的有限维球，计数因而可以完成。
